\chapter{Analyse exploratoire de la population et de la gravité}
\label{chap:exploration}

Une fois la table préparée, le notebook 04 décrit les usagers et les différences de
gravité. Chaque graphique répond à une question et les tableaux conservent les effectifs
qui permettent d'interpréter les pourcentages. Les résultats de ce chapitre sont
descriptifs~; leurs intervalles de variabilité seront examinés au chapitre suivant.

\section{Population analysée et dénominateurs}
Le chargement parcourt le fichier préparé par blocs et conserve uniquement 2019-2022.
Ces années correspondent à la période d'apprentissage. Les comparaisons de gravité de
ce chapitre ne portent pas sur 2023 ou 2024, même si ces années ont participé aux audits.

La population comprend 467\,348 usagers et 207\,057 accidents. La gravité est connue
pour 467\,079 usagers et inconnue pour 269. On compte 83\,081 décès ou hospitalisations,
soit 17,79\,\% des gravités connues. Ce pourcentage est une proportion par usager~;
un accident impliquant plusieurs personnes contribue plusieurs observations.

Pour chaque groupe, le dénominateur est le nombre d'usagers de ce groupe dont la cible
est connue. Par exemple, un cas positif, un cas négatif et une gravité inconnue donnent
une proportion de $1/2$, et non de $1/3$. Un groupe sans gravité connue ne reçoit pas
une proportion de zéro~: cette proportion n'est pas définie. Les effectifs de cibles
inconnues restent présentés dans les sorties du notebook.

\section{Effectifs et profils des usagers}
\subsection{Les volumes enregistrés varient selon l'année}
Le tableau~\ref{tab:population-train} sépare les effectifs d'usagers et d'accidents.
La proportion annuelle se calcule sur les cibles connues de l'année. La proportion
globale résulte de la somme des numérateurs et dénominateurs~; elle n'est pas la moyenne
non pondérée des quatre pourcentages annuels.
\begin{table}[htbp]
\centering\small
\renewcommand{\arraystretch}{1.18}
\caption{Population et proportion de décès ou hospitalisations, 2019-2022.}\label{tab:population-train}
\begin{tabularx}{\textwidth}{@{}lLLLLL@{}}
\toprule
Année & Usagers & Accidents & Gravités connues & Cas positifs & Proportion (\%) \\
\midrule
2019 & 126\,352 & 56\,016 & 126\,352 & 22\,853 & 18,09 \\
2020 & 99\,169 & 45\,121 & 99\,169 & 18\,141 & 18,29 \\
2021 & 122\,189 & 53\,540 & 122\,135 & 20\,796 & 17,03 \\
2022 & 119\,638 & 52\,380 & 119\,423 & 21\,291 & 17,83 \\
\bottomrule
\end{tabularx}
\sourcecalcul{notebook 04 ; dénominateur : gravités connues de chaque année}
\end{table}

\figurerapport{04_01_effectifs_annuels.png}{Évolution des usagers enregistrés et des accidents distincts. Les deux unités utilisent des axes séparés.}{fig:volumes}

La baisse des volumes en 2020 est visible dans les deux unités. Ces courbes ne séparent
pas les effets de l'exposition routière, des pratiques d'enregistrement et des autres
changements de contexte. Les seules données représentées ne suffisent pas à attribuer
la variation à une cause précise ni à en déduire une évolution du risque par trajet.
\FloatBarrier

\subsection{Âge approximatif et catégorie d'usager}
Parmi les 461\,530 âges préparés connus, la médiane est de 35 ans, les quartiles de
23 et 52 ans et la moyenne d'environ 38,5 ans. L'histogramme utilise des classes de
cinq ans. Il décrit la distribution des personnes enregistrées dont l'âge est connu,
sans représenter la population générale française.
\figurerapport[0.90]{04_02_distribution_age.png}{Distribution des âges approximatifs connus. Les 5 818 âges manquants sont comptés séparément.}{fig:ages}

Les conducteurs représentent environ 74,5\,\% des usagers, les passagers 17,7\,\% et
les piétons 7,8\,\%. Ces parts décrivent la composition des accidents enregistrés.
Une catégorie plus nombreuse n'est pas nécessairement celle dont la proportion de
décès ou d'hospitalisations est la plus élevée.
\figurerapport[0.88]{04_03_categories_usagers.png}{Répartition des catégories d'usagers dans la population analysée.}{fig:catu-pop}
\FloatBarrier

\section{Répartition des catégories de gravité}
Les quatre modalités sont conservées pour distinguer ce que le regroupement binaire
réunit. On observe 196\,270 indemnes, 11\,996 tués, 71\,085 blessés hospitalisés et
187\,728 blessés légers, auxquels s'ajoutent les 269 gravités inconnues.
Les pourcentages de la figure~\ref{fig:gravite} portent sur tous les usagers, y compris
les inconnus~; le taux de la cible binaire porte, lui, sur les seules gravités connues.
\figurerapport[0.90]{04_04_repartition_gravite.png}{Répartition des quatre modalités originales et des gravités inconnues, 2019-2022.}{fig:gravite}

Les décès seuls représentent 2,57\,\% des gravités connues, contre 17,79\,\% pour décès
ou hospitalisation. Cette différence de fréquence justifie de lire les deux définitions
séparément. Une méthode prédictive sera comparée à une référence simple et évaluée
au-delà de sa seule proportion globale de réponses correctes.
\FloatBarrier

\section{Comparer les proportions selon les caractéristiques}
\subsection{Âge}
La proportion est de 20,40\,\% chez les 0-17 ans, descend à 13,91\,\% chez les 35-44 ans,
puis atteint 33,26\,\% chez les 75 ans ou plus. Cette forme ne correspond pas à une
progression linéaire unique de la proportion avec l'âge. Les groupes ont aussi des
compositions différentes en matière de déplacement et de contexte.
\figurerapport{04_05_gravite_selon_age.png}{Décès ou hospitalisation selon la tranche d'âge. $n$ désigne les gravités connues ; le trait représente la proportion globale de 17,8\,\%.}{fig:age-gravite}
\begin{table}[htbp]
\centering\small
\renewcommand{\arraystretch}{1.18}
\caption{Gravité selon l'âge, 2019-2022.}\label{tab:gravite-age}
\begin{tabularx}{\textwidth}{@{}>{\hsize=1.65\hsize\linewidth=\hsize}L>{\hsize=.8\hsize\linewidth=\hsize}L>{\hsize=.75\hsize\linewidth=\hsize}L>{\hsize=.8\hsize\linewidth=\hsize}L@{}}
\toprule
Groupe & Gravités connues & Cas positifs & Proportion (\%) \\
\midrule
0-17 ans & 41\,349 & 8\,437 & 20,40 \\
18-24 ans & 87\,449 & 15\,629 & 17,87 \\
25-34 ans & 97\,065 & 13\,814 & 14,23 \\
35-44 ans & 73\,461 & 10\,222 & 13,91 \\
45-54 ans & 64\,377 & 10\,560 & 16,40 \\
55-64 ans & 47\,758 & 9\,752 & 20,42 \\
65-74 ans & 27\,652 & 7\,200 & 26,04 \\
75 ans ou plus & 22\,418 & 7\,457 & 33,26 \\
Âge inconnu & 5\,550 & 10 & 0,18 \\
\bottomrule
\end{tabularx}
\sourcecalcul{notebook 04 ; les gravités inconnues sont exclues du dénominateur}
\end{table}


La tranche 35-44 ans est retenue comme référence descriptive pour la comparaison
avec les 75 ans ou plus. Ce choix est intervenu après l'exploration~: il ne s'agit pas
d'un contraste confirmé à partir d'une hypothèse annoncée avant observation.
\FloatBarrier

\subsection{Catégorie d'usager}
Les proportions sont proches pour les conducteurs et les passagers, autour de 16\,\%,
et plus élevées pour les piétons, à 32,60\,\%. Le tableau~\ref{tab:gravite-categorie}
montre que cette différence ne s'explique pas par la simple taille des groupes~:
les piétons sont moins nombreux mais leur proportion positive est plus forte.
\figurerapport{04_06_gravite_selon_categorie.png}{Proportions selon la catégorie d'usager ; les effectifs indiqués excluent les gravités inconnues.}{fig:catu-gravite}
\begin{table}[htbp]
\centering\small
\renewcommand{\arraystretch}{1.18}
\caption{Gravité selon la catégorie d'usager, 2019-2022.}\label{tab:gravite-categorie}
\begin{tabularx}{\textwidth}{@{}>{\hsize=1.65\hsize\linewidth=\hsize}L>{\hsize=.8\hsize\linewidth=\hsize}L>{\hsize=.75\hsize\linewidth=\hsize}L>{\hsize=.8\hsize\linewidth=\hsize}L@{}}
\toprule
Groupe & Gravités connues & Cas positifs & Proportion (\%) \\
\midrule
Conducteur & 347\,962 & 57\,674 & 16,57 \\
Passager & 82\,591 & 13\,498 & 16,34 \\
Piéton & 36\,526 & 11\,909 & 32,60 \\
\bottomrule
\end{tabularx}
\sourcecalcul{notebook 04 ; les gravités inconnues sont exclues du dénominateur}
\end{table}


Cette comparaison n'isole pas un effet causal du statut de piéton. Elle peut refléter
des différences d'âge, de contexte et de véhicules impliqués. Pour les piétons, le
véhicule associé est celui qui les a heurtés~; une future analyse par catégorie de
véhicule devra conserver cette distinction.
\FloatBarrier

\subsection{Lumière}
La proportion est de 16,94\,\% en plein jour et de 28,79\,\% de nuit sans éclairage
public. Les autres modalités ont des niveaux différents~: 20,14\,\% à l'aube ou au
crépuscule, 19,58\,\% de nuit avec éclairage public éteint et 13,45\,\% avec éclairage
public allumé. Regrouper toutes les nuits masquerait ces différences.
\figurerapport{04_07_gravite_selon_lumiere.png}{Proportions selon la lumière. La modalité non renseignée figure dans le tableau, mais pas dans les barres des modalités connues.}{fig:lum-gravite}
\begin{table}[htbp]
\centering\small
\renewcommand{\arraystretch}{1.18}
\caption{Gravité selon la lumière, 2019-2022.}\label{tab:gravite-lumiere}
\begin{tabularx}{\textwidth}{@{}>{\hsize=1.65\hsize\linewidth=\hsize}L>{\hsize=.8\hsize\linewidth=\hsize}L>{\hsize=.75\hsize\linewidth=\hsize}L>{\hsize=.8\hsize\linewidth=\hsize}L@{}}
\toprule
Groupe & Gravités connues & Cas positifs & Proportion (\%) \\
\midrule
Plein jour & 316\,041 & 53\,546 & 16,94 \\
Crépuscule ou aube & 30\,196 & 6\,082 & 20,14 \\
Nuit sans éclairage public & 45\,404 & 13\,072 & 28,79 \\
Nuit, éclairage public éteint & 3\,872 & 758 & 19,58 \\
Nuit, éclairage public allumé & 71\,557 & 9\,622 & 13,45 \\
Non renseigné & 9 & 1 & 11,11 \\
\bottomrule
\end{tabularx}
\sourcecalcul{notebook 04 ; les gravités inconnues sont exclues du dénominateur}
\end{table}


La faible proportion observée lorsque l'éclairage public est allumé ne mesure pas
l'effet propre d'un éclairage~: les types de routes, l'agglomération et les profils
d'usagers peuvent différer. La modalité non renseignée ne compte que neuf gravités
connues~; son pourcentage ne justifie pas une interprétation substantielle.
\FloatBarrier

\subsection{Agglomération}
La proportion atteint 25,90\,\% hors agglomération, contre 12,89\,\% en agglomération.
Les effectifs restent élevés dans les deux groupes. Cette différence décrit le contexte
des accidents enregistrés~; elle ne fournit ni un risque par kilomètre ni la mesure
d'un effet isolé de la vitesse, qui n'est pas directement observée par \code{vma}.
\figurerapport{04_08_gravite_selon_agglomeration.png}{Proportions de décès ou hospitalisations en et hors agglomération, 2019-2022.}{fig:agg-gravite}
\begin{table}[htbp]
\centering\small
\renewcommand{\arraystretch}{1.18}
\caption{Gravité selon l'agglomération, 2019-2022.}\label{tab:gravite-agglomeration}
\begin{tabularx}{\textwidth}{@{}>{\hsize=1.65\hsize\linewidth=\hsize}L>{\hsize=.8\hsize\linewidth=\hsize}L>{\hsize=.75\hsize\linewidth=\hsize}L>{\hsize=.8\hsize\linewidth=\hsize}L@{}}
\toprule
Groupe & Gravités connues & Cas positifs & Proportion (\%) \\
\midrule
Hors agglomération & 175\,858 & 45\,555 & 25,90 \\
En agglomération & 291\,221 & 37\,526 & 12,89 \\
\bottomrule
\end{tabularx}
\sourcecalcul{notebook 04 ; les gravités inconnues sont exclues du dénominateur}
\end{table}

\FloatBarrier

\section{Valeurs manquantes et lecture annuelle}
Les 5\,818 âges manquants comprennent 5\,550 usagers à gravité connue, dont seulement
10 décès ou hospitalisations, soit 0,18\,\%. Les 268 autres ont une gravité inconnue.
L'absence d'âge est donc associée à une répartition particulière de la gravité. Ce
constat invite à examiner le mécanisme de collecte et la stabilité temporelle d'un
éventuel indicateur d'âge manquant. Il ne permet pas d'affirmer que tous ces usagers
sont en fuite ni de les supprimer par principe.

Les tableaux annuels du notebook reprennent les comparaisons par catégorie d'usager
et agglomération. Leur orientation demeure la même chaque année, mais leurs niveaux
varient. Le chapitre suivant prolonge cette vérification pour les quatre contrastes.
Une persistance annuelle ne constitue pas un ajustement sur les autres caractéristiques.

\section{Ce que l'exploration permet de retenir}
L'âge, la catégorie d'usager, la lumière et l'agglomération présentent des écarts
descriptifs qui méritent d'être quantifiés. Les comparaisons gardent les valeurs
manquantes visibles et précisent leurs dénominateurs. Elles ne recherchent pas toutes
les relations possibles entre les 119 colonnes.

Les codes des catégories ne sont pas traités comme des nombres ordonnés. Les résultats
présentés sont des distributions et des proportions conditionnelles~; aucune matrice
de corrélations de Pearson sur les codes bruts n'est utilisée comme preuve d'association.
Des analyses sur d'autres variables ou avec ajustement simultané resteraient des
travaux supplémentaires, à définir en fonction d'une question précise.

\bilan{L'exploration fournit un portrait de la population et quatre contrastes motivés.
Elle établit leur sens et leur ampleur descriptive. Elle ne mesure pas encore leur
variabilité sous un processus de répétition et ne sépare pas les contributions des
différentes caractéristiques.}
